% =============================================================================
% demo-minimal.tex — Demo document for the swarmmin theme
%
% Compile:  pdflatex demo-minimal.tex   (run twice for TOC)
%      or:  lualatex demo-minimal.tex
%      or:  xelatex  demo-minimal.tex
%
% This demo uses 5 lazy-loaded packages (hyperref, amsmath, enumitem, booktabs, listings)
% to demonstrate the theme's API. Everything else is zero-overhead.
% No shell-escape required. No Pygments. No external dependencies.
% =============================================================================

\documentclass[11pt, a4paper]{article}

\usepackage{swarmmin}

% ── Lazy-load only what this demo actually needs ─────────────────────────────
\usehyperref     % for clickable TOC and cross-references
\usemath         % for theorems and equations (sections 4-5)
\useenumitem     % for styled lists (section 3)
\usebooktabs     % for proper table rules (section 4)
\uselistings     % for code blocks (section 5)

% ── Metadata ─────────────────────────────────────────────────────────────────

\title{SwarmMin Theme Demo}
\subtitle{Ultra-minimal, ultra-fast --- zero packages by default}
\author{LaTeX Helper Swarm Project}
\date{May 15, 2026}

\begin{document}

% ══════════════════════════════════════════════════════════════════════════════
% Title Page
% ══════════════════════════════════════════════════════════════════════════════

\swarmtitlepage

% ══════════════════════════════════════════════════════════════════════════════
% Table of Contents
% ══════════════════════════════════════════════════════════════════════════════

\tableofcontents
\newpage

% ══════════════════════════════════════════════════════════════════════════════
% Section 1: Introduction
% ══════════════════════════════════════════════════════════════════════════════

\section{Introduction}

The \code{swarmmin} theme is the most minimal member of the LaTeX Helper Swarm
theme family. It provides \hltext{API-compatible} commands with
\code{swarmbeauty} while loading \emphtext{zero packages} by default. Every
feature is opt-in via lazy \code{\textbackslash use*} commands: call
\code{\textbackslash usegraphics} for images, \code{\textbackslash usemath}
for equations, \code{\textbackslash uselistings} for code blocks, and so on.
Only load what you actually need.

The design philosophy is radical minimalism: no custom colors, no custom
layouts, no headers or footers, no TOC styling, no section decorations. Block
environments are plain bold-text markers. The title page uses the built-in
\code{titlepage} environment. Table rules lazy-load \code{booktabs} on first
use. The result is the fastest possible compilation while maintaining full API
compatibility with the beautiful and performance themes.

\begin{tipblock}
  Works with pdfLaTeX, XeLaTeX, and LuaLaTeX. No \code{-{}-shell-escape}
  flag is required. No Pygments, no external dependencies, ever.
\end{tipblock}

% ══════════════════════════════════════════════════════════════════════════════
% Section 2: Block Environments
% ══════════════════════════════════════════════════════════════════════════════

\section{Block Environments}

All five block types are plain bold-text markers with zero overhead. No
tcolorbox, no minipage, no colors, no backgrounds --- just a bold label
followed by normal paragraph text.

\begin{noteblock}
  This is a note block. Use it for general observations, additional context,
  or supplementary information that is relevant but not critical to the main
  argument.
\end{noteblock}

\begin{tipblock}
  This is a tip block. Use it to highlight best practices, shortcuts, or
  recommendations that can improve the reader's workflow or understanding.
\end{tipblock}

\begin{warningblock}
  This is a warning block. Use it to alert the reader about potential pitfalls,
  common mistakes, or deprecated features that they should be aware of.
\end{warningblock}

\begin{dangerblock}
  This is an important notice block. Use it for critical information that the
  reader must not overlook --- breaking changes, security issues, or required
  actions.
\end{dangerblock}

\begin{exampleblock}
  This is an example block. Use it for illustrative content, worked examples,
  or step-by-step demonstrations.
\end{exampleblock}

% ══════════════════════════════════════════════════════════════════════════════
% Section 3: Typography and Text
% ══════════════════════════════════════════════════════════════════════════════

\section{Typography and Text}

Inline code: \code{some\_function()}.
Escaped code with special characters: \codeesc{text\_with\_underscores}.
Highlighted text: \hltext{important concept}.
Emphasized text with box: \emphtext{key idea}.

\subsection{Lists}

\begin{itemize}[itemsep=4pt, leftmargin=1.2em]
  \item First item in an unordered list
  \item Second item with \hltext{emphasis}
  \item Third item with nested content:
    \begin{itemize}[itemsep=2pt, leftmargin=1.2em]
      \item Nested item A
      \item Nested item B
    \end{itemize}
\end{itemize}

\begin{enumerate}[itemsep=4pt, leftmargin=1.5em]
  \item Step one: load the minimal theme
  \item Step two: call \code{\textbackslash use*} commands as needed
  \item Step three: compile with any LaTeX engine
\end{enumerate}

\colorrule

% ══════════════════════════════════════════════════════════════════════════════
% Section 4: Tables
% ══════════════════════════════════════════════════════════════════════════════

\section{Tables}

Table rules (\code{\textbackslash swarmtoprule}, \code{\textbackslash swarmmidrule},
\code{\textbackslash swarmbottomrule}) lazy-load \code{booktabs} on first use.

\begin{table}[h]
  \centering
  \caption{Swarm theme comparison}
  \label{tab:comparison}
  \begin{tabular}{l c c c}
    \swarmtoprule
    \textbf{Theme} & \textbf{Packages at load} & \textbf{Features} & \textbf{Engine} \\
    \swarmmidrule
    swarmbeauty & 20+                & Full styling, TikZ, minted & LuaLaTeX only \\
    swarmperf   & 13                 & Colors, headers, listings  & Any           \\
    swarmmin    & 0 (lazy loading)   & Bare API, opt-in packages  & Any           \\
    \swarmbottomrule
  \end{tabular}
\end{table}

See Table~\ref{tab:comparison} for a feature overview.

% ══════════════════════════════════════════════════════════════════════════════
% Section 5: Code Listings
% ══════════════════════════════════════════════════════════════════════════════

\section{Code Listings}

Code blocks use the \code{listings} package (loaded via
\code{\textbackslash uselistings}). Zero external dependencies --- no Pygments,
no shell-escape required.

\begin{lstlisting}[language=Python,
  title={\sffamily\small Fibonacci sequence generator}]
def fibonacci(n: int) -> list[int]:
    """Generate the first n Fibonacci numbers."""
    seq = [0, 1]
    for i in range(2, n):
        seq.append(seq[-1] + seq[-2])
    return seq[:n]

if __name__ == "__main__":
    result = fibonacci(10)
    print(f"First 10: {result}")
\end{lstlisting}

\begin{lstlisting}[language=TeX]
\begin{equation}
  E = mc^2
\end{equation}
\end{lstlisting}

% ══════════════════════════════════════════════════════════════════════════════
% Section 6: Mathematics
% ══════════════════════════════════════════════════════════════════════════════

\section{Mathematics}

Theorem environments use plain text formatting with section-based numbering.
API-compatible with swarmbeauty: \code{\{name\}\{label\}} (both mandatory).

\begin{theorem}{Euler's Identity}{euler}
  For any real number $x$, the following identity holds:
  \[
    e^{i\pi} + 1 = 0
  \]
\end{theorem}

\begin{definition}{Metric Space}{metric}
  A \emph{metric space} is a pair $(X, d)$ where $X$ is a set and
  $d\colon X \times X \to \mathbb{R}_{\geq 0}$ satisfies:
  \begin{enumerate}[itemsep=2pt]
    \item $d(x, y) = 0 \iff x = y$ \quad (identity of indiscernibles)
    \item $d(x, y) = d(y, x)$ \quad (symmetry)
    \item $d(x, z) \leq d(x, y) + d(y, z)$ \quad (triangle inequality)
  \end{enumerate}
\end{definition}

\begin{lemma}{Nested Intervals}{nested}
  Let $\{[a_n, b_n]\}_{n=1}^{\infty}$ be a sequence of closed, bounded
  intervals such that $[a_{n+1}, b_{n+1}] \subseteq [a_n, b_n]$ for all $n$.
  Then $\bigcap_{n=1}^{\infty} [a_n, b_n] \neq \emptyset$.
\end{lemma}

\subsection{Inline and Display Math}

The quadratic formula: $x = \frac{-b \pm \sqrt{b^2 - 4ac}}{2a}$

A more complex display equation:
\[
  \int_{-\infty}^{\infty} e^{-x^2} \, dx = \sqrt{\pi}
\]

\end{document}
